\chapter{Discussion and Conclusions}
\label{chap:discussion-outlook}

This thesis set out to test whether a differentiable product-GP representation
could make the mapping-QCLE excess term usable on a moving trajectory cloud.
The answer for the construction studied here is no. Accurate interpolation of
density values on a fixed cloud does not imply accurate ambient derivatives,
preservation of the physical SEO image, stable signed estimators, or
conservative dynamics. The contribution is therefore a precisely defined
formulation, a sequence of physically interpretable tests that expose its
limitations, and concrete requirements for a successor. It is not a validated
GP-QCLE solver and it is not evidence that MIDPOINT improves PBME.

\section{Application-specific contribution}
\label{sec:discussion-contribution}

Values on a lower-dimensional manifold do not generally identify ambient
normal derivatives; that fact is not original to this thesis. The original
application-specific contribution is narrower: constructing a differentiable
product-GP excess-source prototype on moving PBME/MInt clouds, testing it
against an analytic manufactured operator and the exact two-state SEO image,
and connecting its observed derivative, projection, conservation, and
signed-weight failures to the propagated dynamics without claiming a unique
cause
\cite{Kim2008JCP,Kelly2012JCP,Raissi2017LinearGP,
Solak2002DerivativeGP,CockayneOatesSullivanGirolami2019}.

\section{Physical interpretation of the failure}
\label{sec:discussion-evidence}

The results in Chapter~\ref{chap:results} reveal a connected, but not uniquely
causal, failure pathway. A product kernel can reconstruct the sampled density
values while leaving ambient derivatives underdetermined away from the cloud.
The excess operator depends precisely on those derivatives. Once its source
is applied without enforcing the four-field SEO structure, the fitted density
can develop large components outside the physical image. The focused-sampling
ratio then combines sign cancellation with very large individual label ratios,
which magnifies cloud-to-cloud variation. The resulting raw normalization,
trace, and energy drift show that this is not merely a cosmetic change in
density shape.

These observations establish compatibility among the failure modes, not a
proof that one is the sole cause of the others. The manufactured operator
study, SEO projection diagnostic, signed-label analysis, and conservation
test deliberately isolate different links in the chain. Their agreement
supports the method-level conclusion without overstating causal identification.
The controlled TDSE and grid-QCLE studies then ensure that the absence of a
systematic physical improvement is not an artefact of an uncontrolled
reference calculation.

\section{Claim boundary}
\label{sec:discussion-limitations}

This thesis does not establish a validated discretization of the complete
mapping-basis QCLE, deterministic cloud convergence, strong uncertainty
calibration from four seeds, chemical accuracy outside the one-dimensional
two-state benchmark, multidimensional scalability, or systematic improvement
over PBME. TDSE and grid-QCLE have different epistemic roles: TDSE is
model-exact within numerical controls, whereas grid QCLE solves an approximate
quantum--classical equation numerically.

The production regularization \(0.05\), pilot \(0.01\), and
manufactured-test \(10^{-6}\) policies are all tested and reported rather than
conflated. The tail study demonstrates extreme signed-weight instability, but
its negligible-mass rule does not permit a nontrivial point-removal plateau,
so the instability is not assigned uniquely to small initial labels.

\section{Requirements for a successor}
\label{sec:discussion-successor}

A credible successor should regress the four real bath-dependent fields of the
two-state SEO image so that Hermiticity and projection are enforced by
construction. Its discrete excess source should preserve normalization and
the relevant trace and energy identities before any renormalization. The
regression should incorporate derivative information or structural constraints
capable of controlling directions normal to the sampled manifold. A future
success claim requires decreasing manufactured-operator errors under a
controlled refinement, time-step changes resolvable above seed variation,
independent-cloud uncertainty bands, raw conservation, and reproducibly
smaller PBME-matched errors against controlled references where the excess
source is appreciable
\cite{TraskBochevPerego2020,BonetLok1999,Monaghan2005}.

\section{Numerical scope and transferability}
\label{sec:discussion-numerical-scope}

The numerical conclusions are tied to the parameter ranges, independent seed
sets, domains, and discretizations stated in Chapter~\ref{chap:results} and
Appendix~\ref{app:reference-and-reproducibility}. Repeating those calculations
tests this particular construction on the same benchmark; it does not by
itself establish transferability to additional electronic states, nuclear
dimensions, coupling regimes, or sampling measures. Such transfer would
require renewed manufactured-operator, projection, conservation, and
reference comparisons rather than extrapolation from the present negative
result.

\section{Conclusion}
\label{sec:discussion-conclusions}

The combined numerical evidence supports a controlled negative result.
Accurate same-cloud value reconstruction does not guarantee an accurate
derivative-dependent excess operator, preservation of the SEO image, raw
conservation, or seed-stable dynamics. The tested MIDPOINT branch fails the
joint reliability criteria and does not demonstrate systematic improvement
over PBME. This conclusion applies to the tested discretization, not to the
continuum QCLE excess term.

Accordingly, the thesis contributes a precisely scoped formulation, a
physics-based failure analysis, and evidence-backed redesign criteria. It
does not claim that the tested construction is reliable, deterministically
cloud-converged, or improved.
